\documentclass[12pt,a4paper]{article}
\usepackage[UTF8]{ctex}
\usepackage{amsmath,amssymb,amsthm}
\usepackage{geometry}
\usepackage{bm}
\usepackage{hyperref}
\usepackage{tikz}
\usetikzlibrary{arrows,arrows.meta,positioning,calc,shapes.geometric}

\geometry{margin=2.5cm}

\title{软接触（Soft Contact）的详细理解}
\author{第12章抓取与操作补充说明}
\date{}

\newtheorem{theorem}{定理}
\newtheorem{definition}{定义}
\newtheorem{example}{例子}

\begin{document}

\maketitle

\section{概述}

\textbf{软接触}（soft contact），也称为\textbf{软指接触}（soft-finger contact），是接触运动学中的三种点接触模型之一。它是三种模型中约束最严格的一种，能够模拟由于接触面积非零而产生的额外约束。

\section{三种点接触模型}

在接触运动学中，为了在不显式建模力的情况下近似摩擦效果，我们定义了三种纯运动学点接触模型：

\subsection{1. 无摩擦点接触（Frictionless Point Contact）}

\begin{itemize}
\item \textbf{约束}：只强制执行滚动-滑动约束(12.5)
  \begin{equation}
  F^T (V_A - V_B) \geq 0
  \end{equation}
\item \textbf{允许的运动}：
  \begin{itemize}
  \item 分离（Breaking）：$F^T (V_A - V_B) > 0$
  \item 滑动（Sliding）：$F^T (V_A - V_B) = 0$，但$\dot{p}_A \neq \dot{p}_B$
  \item 滚动（Rolling）：$F^T (V_A - V_B) = 0$，且$\dot{p}_A = \dot{p}_B$
  \end{itemize}
\item \textbf{特点}：约束最少，允许接触点处有相对滑动
\end{itemize}

\subsection{2. 有摩擦点接触（Point Contact with Friction）}

\begin{itemize}
\item \textbf{约束}：强制执行滚动约束(12.8)
  \begin{equation}
  \dot{p}_A = v_A + [\omega_A]p_A = v_B + [\omega_B]p_B = \dot{p}_B
  \end{equation}
\item \textbf{允许的运动}：
  \begin{itemize}
  \item 分离（Breaking）：$F^T (V_A - V_B) > 0$
  \item 滚动（Rolling）：$\dot{p}_A = \dot{p}_B$（接触点处无相对运动）
  \item \textbf{不允许滑动}
  \end{itemize}
\item \textbf{特点}：隐式建模了足以防止接触处滑动的摩擦力
\item \textbf{物理意义}：摩擦力足够大，使得接触点处不能有切向相对运动
\end{itemize}

\subsection{3. 软接触（Soft Contact / Soft-Finger Contact）}

\begin{itemize}
\item \textbf{约束}：
  \begin{enumerate}
  \item 滚动约束(12.8)：$\dot{p}_A = \dot{p}_B$
  \item \textbf{额外的约束}：两个物体不能绕接触法线轴相对旋转
  \end{enumerate}
\item \textbf{允许的运动}：
  \begin{itemize}
  \item 分离（Breaking）：$F^T (V_A - V_B) > 0$
  \item 滚动（Rolling）：$\dot{p}_A = \dot{p}_B$，且绕法线轴无相对旋转
  \item \textbf{不允许滑动，也不允许绕法线轴旋转}
  \end{itemize}
\item \textbf{特点}：约束最严格
\end{itemize}

\section{软接触的详细分析}

\subsection{额外约束的数学表达}

软接触除了滚动约束外，还要求两个物体不能绕接触法线轴相对旋转。

设接触法线为$\hat{n}$，则额外约束为：
\begin{equation}
\hat{n}^T (\omega_A - \omega_B) = 0
\label{eq:soft_contact_constraint}
\end{equation}

这意味着：两个物体的角速度在接触法线方向上的分量必须相等。

\subsection{为什么叫"软"接触？}

\textbf{"软"的含义}：
\begin{itemize}
\item 不是指材料软，而是指接触时物体会发生\textbf{变形}
\item 由于变形，接触不是理想点接触，而是有\textbf{非零接触面积}
\item 非零接触面积导致可以传递\textbf{摩擦力矩}（friction moment）
\item 摩擦力矩抵抗绕接触法线轴的旋转
\end{itemize}

\subsection{物理机制}

\begin{enumerate}
\item \textbf{接触变形}：
  \begin{itemize}
  \item 当两个物体接触时，如果材料较软或接触力较大，会发生变形
  \item 变形导致接触不是点接触，而是有接触面积
  \end{itemize}

\item \textbf{摩擦力矩}：
  \begin{itemize}
  \item 由于有接触面积，摩擦力不仅产生在切向，还产生\textbf{绕法线轴的力矩}
  \item 这个力矩抵抗两个物体绕接触法线轴的相对旋转
  \end{itemize}

\item \textbf{运动学约束}：
  \begin{itemize}
  \item 从运动学角度，我们用一个额外的约束来模拟这个效果
  \item 约束：$\hat{n}^T (\omega_A - \omega_B) = 0$
  \end{itemize}
\end{enumerate}

\section{三种接触模型的对比}

\subsection{约束数量}

\begin{center}
\begin{tabular}{|c|c|c|}
\hline
\textbf{接触模型} & \textbf{约束数量} & \textbf{约束内容} \\
\hline
无摩擦点接触 & 1个 & 滚动-滑动约束 \\
\hline
有摩擦点接触 & 3个 & 滚动约束（接触点处无相对运动） \\
\hline
软接触 & 4个 & 滚动约束 + 绕法线轴无相对旋转 \\
\hline
\end{tabular}
\end{center}

\subsection{允许的自由度}

对于3D情况，两个物体在接触点处的相对运动有6个自由度：
\begin{itemize}
\item 3个平移自由度：$v_A - v_B$
\item 3个旋转自由度：$\omega_A - \omega_B$
\end{itemize}

\begin{center}
\begin{tabular}{|c|c|}
\hline
\textbf{接触模型} & \textbf{允许的相对运动} \\
\hline
无摩擦点接触 & 法线方向分离，切向滑动，任意旋转 \\
\hline
有摩擦点接触 & 法线方向分离，无滑动，任意旋转 \\
\hline
软接触 & 法线方向分离，无滑动，无绕法线轴旋转 \\
\hline
\end{tabular}
\end{center}

\subsection{约束方程总结}

\textbf{无摩擦点接触}：
\begin{equation}
F^T (V_A - V_B) \geq 0
\end{equation}
其中$F = ([p]\hat{n}, \hat{n})$是接触wrench。

\textbf{有摩擦点接触}：
\begin{align}
F^T (V_A - V_B) &\geq 0 \quad \text{（不可穿透）} \\
\dot{p}_A &= \dot{p}_B \quad \text{（滚动约束，3个标量方程）}
\end{align}

\textbf{软接触}：
\begin{align}
F^T (V_A - V_B) &\geq 0 \quad \text{（不可穿透）} \\
\dot{p}_A &= \dot{p}_B \quad \text{（滚动约束，3个标量方程）} \\
\hat{n}^T (\omega_A - \omega_B) &= 0 \quad \text{（绕法线轴无相对旋转，1个标量方程）}
\end{align}

\section{平面情况下的特殊性质}

\subsection{平面问题的特点}

对于\textbf{平面问题}（planar problems）：
\begin{itemize}
\item 物体限制在平面内运动
\item 只有绕垂直于平面的轴的旋转（$\omega_z$）
\item 接触法线也在平面内
\end{itemize}

\subsection{有摩擦点接触与软接触的等价性}

\textbf{重要结论}：对于平面问题，\textbf{有摩擦点接触和软接触是等价的}。

\textbf{原因}：
\begin{itemize}
\item 在平面问题中，旋转只有$\omega_z$分量
\item 接触法线$\hat{n}$在平面内，垂直于$z$轴
\item 因此：$\hat{n}^T \omega = 0$（法线垂直于旋转轴）
\item 约束$\hat{n}^T (\omega_A - \omega_B) = 0$自动满足
\item 软接触的额外约束在平面情况下是冗余的
\end{itemize}

\textbf{数学证明}：

在平面问题中：
\begin{itemize}
\item 角速度：$\omega_A = [0, 0, \omega_{Az}]^T$，$\omega_B = [0, 0, \omega_{Bz}]^T$
\item 接触法线：$\hat{n} = [n_x, n_y, 0]^T$（在$xy$平面内）
\end{itemize}

因此：
\begin{align}
\hat{n}^T (\omega_A - \omega_B) &= [n_x, n_y, 0] \begin{bmatrix} 0 \\ 0 \\ \omega_{Az} - \omega_{Bz} \end{bmatrix} \\
&= 0
\end{align}

所以约束$\hat{n}^T (\omega_A - \omega_B) = 0$总是成立，是冗余约束。

\section{实际应用}

\subsection{何时使用软接触模型？}

软接触模型适用于：
\begin{itemize}
\item \textbf{软材料}：橡胶、硅胶等容易变形的材料
\item \textbf{大接触力}：即使材料较硬，如果接触力很大，也会产生变形
\item \textbf{需要精确控制旋转}：当需要防止物体绕接触法线轴旋转时
\end{itemize}

\subsection{例子：机器人手指抓取}

\begin{itemize}
\item \textbf{硬手指}：使用有摩擦点接触模型
  \begin{itemize}
  \item 接触面积小，近似点接触
  \item 可以传递切向摩擦力，但不能传递绕法线轴的力矩
  \end{itemize}

\item \textbf{软手指}：使用软接触模型
  \begin{itemize}
  \item 手指材料软，接触时发生变形
  \item 有接触面积，可以传递绕法线轴的摩擦力矩
  \item 可以更好地控制被抓取物体的旋转
  \end{itemize}
\end{itemize}

\section{约束的维度分析}

\subsection{3D情况}

对于3D情况，两个物体的相对运动有6个自由度：

\begin{itemize}
\item \textbf{无摩擦点接触}：
  \begin{itemize}
  \item 1个约束（不可穿透）
  \item 剩余5个自由度
  \end{itemize}

\item \textbf{有摩擦点接触}：
  \begin{itemize}
  \item 1个不可穿透约束 + 3个滚动约束 = 4个约束
  \item 剩余2个自由度（可以绕法线轴旋转，可以分离）
  \end{itemize}

\item \textbf{软接触}：
  \begin{itemize}
  \item 1个不可穿透约束 + 3个滚动约束 + 1个旋转约束 = 5个约束
  \item 剩余1个自由度（只能分离）
  \end{itemize}
\end{itemize}

\subsection{平面情况}

对于平面情况，两个物体的相对运动有3个自由度（$v_x, v_y, \omega_z$）：

\begin{itemize}
\item \textbf{无摩擦点接触}：
  \begin{itemize}
  \item 1个约束
  \item 剩余2个自由度
  \end{itemize}

\item \textbf{有摩擦点接触和软接触}：
  \begin{itemize}
  \item 1个不可穿透约束 + 2个滚动约束 = 3个约束
  \item 剩余0个自由度（完全固定，除非分离）
  \end{itemize}
\end{itemize}

\section{总结}

\subsection{关键点}

\begin{enumerate}
\item \textbf{软接触}是三种点接触模型中约束最严格的一种

\item \textbf{额外约束}：两个物体不能绕接触法线轴相对旋转
  \begin{equation}
  \hat{n}^T (\omega_A - \omega_B) = 0
  \end{equation}

\item \textbf{物理机制}：由于接触变形和非零接触面积，可以传递绕法线轴的摩擦力矩

\item \textbf{平面情况}：对于平面问题，有摩擦点接触和软接触是等价的

\item \textbf{应用}：适用于软材料或需要防止绕法线轴旋转的情况
\end{enumerate}

\subsection{三种模型的层次结构}

\begin{center}
\begin{tikzpicture}[node distance=2cm]
% 定义节点
\node[draw,rectangle] (frictionless) {无摩擦点接触\\1个约束};
\node[draw,rectangle,below of=frictionless] (friction) {有摩擦点接触\\4个约束};
\node[draw,rectangle,below of=friction] (soft) {软接触\\5个约束};

% 连接
\draw[->] (frictionless) -- (friction) node[midway,right] {增加滚动约束};
\draw[->] (friction) -- (soft) node[midway,right] {增加旋转约束};

% 标签
\node[left of=frictionless,xshift=-1cm] {约束最少};
\node[left of=soft,xshift=-1cm] {约束最多};
\end{tikzpicture}
\end{center}

约束严格程度：无摩擦点接触 $<$ 有摩擦点接触 $<$ 软接触

\end{document}

